\documentclass[28pt, a0paper, landscape, margin=3mm,
               innermargin=8mm, blockverticalspace=20mm,
               colspace=15mm, subcolspace=8mm]{tikzposter}

\usepackage[utf8]{inputenc}
\usepackage[T1]{fontenc}
\usepackage[czech]{babel}
\usepackage{amsmath, amssymb}
\usepackage{graphicx}
\usepackage{booktabs}

\graphicspath{{../output/}}

% ------------------------------------------------------------------
% Barevné schéma a styl bloků (ploché obdélníky, tmavě modrá hlavička)
% ------------------------------------------------------------------
\definecolorstyle{MSMStyle}{
    \definecolor{colorOne}{RGB}{17,54,97}   % tmavě modrá - title/block hlavičky
    \definecolor{colorTwo}{RGB}{17,54,97}
    \definecolor{colorThree}{RGB}{245,247,250}
}{
    \colorlet{backgroundcolor}{white}
    \colorlet{framecolor}{colorOne}
    \colorlet{titlefgcolor}{white}
    \colorlet{titlebgcolor}{colorOne}
    \colorlet{blocktitlebgcolor}{colorOne}
    \colorlet{blocktitlefgcolor}{white}
    \colorlet{blockbodybgcolor}{white}
    \colorlet{blockbodyfgcolor}{black}
    \colorlet{notefgcolor}{black}
    \colorlet{notefrcolor}{colorOne}
    \colorlet{notebgcolor}{colorThree}
}

\usecolorstyle{MSMStyle}
\usetheme{Default}
\useblockstyle[roundedcorners=0pt, titlewidthscale=1, bodywidthscale=1]{Basic}
\usetitlestyle[roundedcorners=0pt]{Filled}

\tikzposterlatexaffectionproofoff

\title{\parbox{0.95\linewidth}{\centering Vícerozměrné náhodné veličiny: generování a charakteristiky\\
\large obecného dvourozměrného rozdělení s hustotou $f(x,y) = c\,x\,y$}}
\author{Michael Hladky \and Vít Novotný}
\institute{ZČU FAV -- KMA/MSM}

\begin{document}
\maketitle

\begin{columns}

% ===================== SLOUPEC 1 =====================
\column{0.25}

\block{Úvod}{
Cílem projektu je analyzovat spojitou náhodnou veličinu $(X,Y)$, definovanou
na obecné dvourozměrné oblasti $\Omega$. Byla zvolena veličina $(X,Y)$,
jejíž hustota pravděpodobnosti je dána funkcí s neznámou konstantou $c$:
\begin{equation*}
f(x,y) = c\,x\,y
\end{equation*}
\begin{equation*}
\Omega = \{(x,y) : 1 < x < 2,\ x < y < 2\}
\end{equation*}
}

\block{Technologie a program}{
Pro výpočty a simulace byl použit programovací jazyk Python s knihovnami
NumPy, SymPy, SciPy a Matplotlib. Program \texttt{main.py} je předem
nastaven na zadanou funkci a po spuštění provede všechny potřebné výpočty
(symbolicky pomocí SymPy) a generování náhodných hodnot metodou zamítání.
Zároveň vytvoří grafické výstupy do adresáře \texttt{output/}.
}

\block{Normovací parametr $c$ a hustota pravděpodobnosti}{
Normovací parametr $c$ lze určit z podmínky, že celková pravděpodobnost
v oblasti $\Omega$ je rovna jedné:
\begin{equation*}
\int_{1}^{2}\int_{x}^{2} c\,x\,y \;\mathrm{d}y\,\mathrm{d}x = 1 .
\end{equation*}
Po výpočtu integrálu získáme hodnotu $c = \dfrac{8}{9}$. Hustota
pravděpodobnosti je tedy
\begin{equation*}
f(x,y) = \frac{8}{9}\,x\,y, \qquad (x,y)\in\Omega .
\end{equation*}

\centering
\includegraphics[width=\linewidth]{part1_joint_density.pdf}
}

% ===================== SLOUPEC 2 =====================
\column{0.25}

\block{Marginální funkce hustoty a distribuční funkce}{
Funkce $f_X(x)$ a $f_Y(y)$ jsou určeny integrací hustoty pravděpodobnosti
přes druhou proměnnou:
\begin{equation*}
f_X(x) = \int_{x}^{2} f(x,y)\,\mathrm{d}y = \frac{16}{9}x - \frac{4}{9}x^3
\end{equation*}
\begin{equation*}
f_Y(y) = \int_{1}^{y} f(x,y)\,\mathrm{d}x = \frac{4}{9}y^3 - \frac{4}{9}y
\end{equation*}
pro $1<x<2$, resp. $1<y<2$.

Distribuční funkce $F_X(x)$ a $F_Y(y)$ jsou určeny integrací marginálních
funkcí hustoty:
\begin{equation*}
F_X(x) = -\frac{1}{9}x^4 + \frac{8}{9}x^2 - \frac{7}{9}
\end{equation*}
\begin{equation*}
F_Y(y) = \frac{1}{9}y^4 - \frac{2}{9}y^2 + \frac{1}{9}
\end{equation*}

\centering
\includegraphics[width=\linewidth]{part2_marginal_functions.pdf}
}

\block{Střední hodnoty, rozptyly a kovariance}{
\begin{equation*}
E[X] = \frac{188}{135} \approx 1.3926
\end{equation*}
\begin{equation*}
E[Y] = \frac{232}{135} \approx 1.7185
\end{equation*}
\begin{equation*}
\operatorname{Var}(X) = \frac{1106}{18225} \approx 0.06069
\end{equation*}
\begin{equation*}
\operatorname{Var}(Y) = \frac{851}{18225} \approx 0.04669
\end{equation*}
\begin{equation*}
\operatorname{cov}(X,Y) = \frac{484}{18225} \approx 0.02656
\end{equation*}

Variační matice:
\begin{equation*}
\operatorname{Var}(X,Y) =
\begin{pmatrix} \operatorname{Var}(X) & \operatorname{cov}(X,Y) \\
\operatorname{cov}(X,Y) & \operatorname{Var}(Y) \end{pmatrix}
\end{equation*}
\begin{equation*}
\approx
\begin{pmatrix} 0.06069 & 0.02656 \\ 0.02656 & 0.04669 \end{pmatrix}
\end{equation*}
}

% ===================== SLOUPEC 3 =====================
\column{0.25}

\block{Generování náhodných hodnot a bodové odhady}{
Náhodný výběr byl generován metodou zamítání (\emph{acceptance--rejection}):
kandidáti $(x,y)$ se losují rovnoměrně na obdélníku $[1,2]\times[1,2]$,
podmínka $y \ge x$ zajistí platnost oblasti $\Omega$ a bod je přijat,
pokud $u \le f(x,y)$ pro $u \sim U(0, f_{\max})$, kde
$f_{\max} = f(2,2) = \tfrac{32}{9} \approx 3.556$.

Pro $n = 1000$ vygenerovaných bodů byly z výběru odhadnuty charakteristiky:
\begin{equation*}
\hat E[X] \approx 1.4014, \qquad \hat E[Y] \approx 1.7285
\end{equation*}
\begin{equation*}
\widehat{\operatorname{Var}}(X) \approx 0.06042, \qquad
\widehat{\operatorname{Var}}(Y) \approx 0.04207
\end{equation*}
\begin{equation*}
\widehat{\operatorname{cov}}(X,Y) \approx 0.02391
\end{equation*}

95\,\% intervalové odhady pro střední hodnoty:
\begin{equation*}
E[X] \in (1.3861,\ 1.4166)
\end{equation*}
\begin{equation*}
E[Y] \in (1.7158,\ 1.7412)
\end{equation*}
}

\block{Porovnání s teoretickými funkcemi}{
Histogramy a empirické distribuční funkce sestrojené z vygenerovaných dat
($n=1000$) byly porovnány se spočtenými teoretickými marginálními funkcemi
hustoty a distribučními funkcemi (viz blok výše). Shoda je ve všech
případech velmi dobrá -- histogramy i empirické distribuční funkce
kopírují teoretické křivky, drobné odchylky jsou dány konečným rozsahem
výběru a náhodností generování.

\centering
\includegraphics[width=\linewidth]{part6_comparison_with_theory.pdf}
}

% ===================== SLOUPEC 4 =====================
\column{0.25}

\block{Závislost odhadů na rozsahu výběru}{
Ukázka chování výběrových odhadů pro rostoucí $n$. Grafy ilustrují, že
pro $n \to \infty$ platí $\bar X_n \to E[X]$, $\bar Y_n \to E[Y]$ a šířka
intervalu spolehlivosti klesá úměrně $1/\sqrt{n}$ -- desetinásobné zvětšení
výběru tedy zúží interval přibližně $\sqrt{10}\approx3.16\times$.

\centering
\begin{minipage}{0.5\linewidth}
\centering
\includegraphics[width=\linewidth]{part7a_estimate_x.pdf}
\end{minipage}%
\begin{minipage}{0.5\linewidth}
\centering
\includegraphics[width=\linewidth]{part7b_estimate_y.pdf}
\end{minipage}

\includegraphics[width=0.85\linewidth]{part7c_ci_width.pdf}

\vspace{1em}
\centering
\small
\begin{tabular}{rrrrr}
\toprule
$n$ & $\hat E[X]$ & šířka IS$_X$ & $\hat E[Y]$ & šířka IS$_Y$ \\
\midrule
10      & 1.3713 & 0.4248 & 1.7334 & 0.2310 \\
100     & 1.3704 & 0.0956 & 1.6884 & 0.0857 \\
1\,000   & 1.3915 & 0.0303 & 1.7307 & 0.0264 \\
10\,000  & 1.3924 & 0.0096 & 1.7188 & 0.0085 \\
100\,000 & 1.3933 & 0.0031 & 1.7169 & 0.0027 \\
\bottomrule
\end{tabular}

\normalsize
\vspace{0.5em}
Teoretické hodnoty: $E[X]=1.3926$, $E[Y]=1.7185$.
}

\block{Závěr}{
Simulace potvrdila teoretické výpočty: hustota $f(x,y)=\tfrac{8}{9}xy$
na $\Omega$ byla úspěšně validována metodou zamítání -- bodové odhady
z $n=1000$ vygenerovaných bodů ($\hat E[X]\approx1.4014$,
$\hat E[Y]\approx1.7285$) leží blízko teoretických hodnot
($E[X]=1.3926$, $E[Y]=1.7185$) a jsou obsaženy v~95\,\% intervalech
spolehlivosti. S~rostoucím rozsahem výběru $n$ se odhady systematicky
zpřesňují a šířka intervalů spolehlivosti klesá úměrně $1/\sqrt{n}$,
což odpovídá očekávanému chování podle centrální limitní věty.
}

\end{columns}

\end{document}
